\section{数据集与分布问题}

\subsection{数据质量的三维观察}

Offline RL 的数据集并非黑盒，讲者强调：理解数据的来源、覆盖范围和奖励分布，才能有效指导算法选择。

\begin{itemize}
    \item \textbf{覆盖度}：状态-动作组合是否均匀，是否存在稀疏区域
    \item \textbf{奖励结构}：奖励是否稠密、有噪声、有异常值
    \item \textbf{行为多样性}：是否仅包含一条行为策略，或多种策略的混合
\end{itemize}

\begin{importantbox}{用数据质量驱动算法选择}
当数据覆盖面广、奖励干净时，更激进的 bootstrapping 方法（如 IQL）可以发挥；当数据存在低质量样本或奖励漂移时，优先使用 Conservative 方法或基于 reward 的 sample filtering。
\end{importantbox}

\begin{table}[H]
\centering
\begin{tabular}{p{0.2\textwidth} p{0.25\textwidth} p{0.25\textwidth} p{0.2\textwidth}}
\toprule
数据特征 & 示意说明 & 风险 & 推荐策略 \\
\midrule
高奖励、低覆盖 & 同一轨迹重复出现 & 过拟合、缺乏泛化 & 加强 regularization、引入数据增强 \\
低奖励但多样 & 多个 suboptimal 策略 & Q overestimation on rare good actions & 过滤低奖励样本、使用 expectile distortion \\
混合行为 & 多策略混合而未标注 & 难以定义行为策略 & 赋予 behavior cloning 更高权重的分段训练 \\
\bottomrule
\end{tabular}
\caption{根据数据分布选择 offline RL 策略}
\end{table}

\subsection{本章小结}

Offline RL 的成功依赖于对数据覆盖度、奖励质量和行为多样性的精准理解。按数据特征调整策略（filter/IQL/CQL）比固定算法更稳妥。

\subsection{案例：构建混合数据管线}

项目组经常将 simulator logs、human logs 与 expert trajectories 混合，按照 reward 和 importance sampling 重新加权样本：

\begin{enumerate}
    \item 先计算每条轨迹的 effective sample size。
    \item 按 reward percentile 将数据分成 bin 并赋予不同 sampling weight。
    \item 用 IQL expectile 评估每个 bin 中 Q 的分布，两端重采样以保持稳定。
    \item 若 expectile 跳出预设区间，退回 conservative 模式或过滤掉异常 bin。
\end{enumerate}

\begin{knowledgebox}{混合数据管线的好处}
不同数据源 bias 各异：simulator 更精确但缺乏多样性，human data 有噪声但更自然。混合管线通过 re-weighting 让算法既稳健又具泛化。
\end{knowledgebox}

